\documentclass[11pt]{article}
\usepackage[paperwidth=210mm,paperheight=5000mm,margin=20mm]{geometry}
\usepackage{amsmath,amssymb,xcolor,graphicx}
\usepackage[hypcap=false]{caption}
\usepackage[hidelinks,bookmarksnumbered,bookmarksopen]{hyperref}
\hypersetup{pdftitle={Mechanical properties — Dislocations and strengthening},pdfauthor={Lecturer: Prof. J. Adeyemi},
  pdfcreator={KherveNote}}
\renewcommand\labelitemii{\textbullet}
\renewcommand\labelitemiii{\textbullet}
\renewcommand\labelitemiv{\textbullet}
\renewcommand\labelenumii{\arabic{enumi}.\arabic{enumii}.}
\renewcommand\labelenumiii{\arabic{enumi}.\arabic{enumii}.\arabic{enumiii}.}
\renewcommand\labelenumiv{\arabic{enumi}.\arabic{enumii}.\arabic{enumiii}.\arabic{enumiv}.}
\setlength{\parindent}{0pt}
\setlength{\parskip}{0.6em}
\definecolor{knoteaccent}{HTML}{1A6DD8}
\definecolor{knotemuted}{HTML}{6B6B6B}
\definecolor{knotekey}{HTML}{D96B00}
\newcommand\knotetime[1]{\leavevmode\llap{\color{knotemuted}\scriptsize #1\hspace{1em}}}
\newcommand\knotekey[2][]{\par\noindent#1\colorbox{knotekey!12}{\parbox{\dimexpr\linewidth-2\fboxsep}{%
  \textbf{\color{knotekey}$\star$ Key point.} #2}}\par}
\newcommand\knotequestion[2][]{\par\noindent#1{\color{knoteaccent}\textbf{?}}~\emph{#2}\par}
\newenvironment{knotetranscript}{\par\color{knotemuted}\small}{\par}
\newcommand\knotesummary[1]{\par\noindent\fcolorbox{knoteaccent}{knoteaccent!6}{%
  \parbox{\dimexpr\linewidth-2\fboxsep-2\fboxrule}{\textbf{Summary}\par #1}}\par}
\makeatletter
\newdimen\knote@limit \knote@limit=4950mm
\newbox\knote@page \newbox\knote@chunk
\def\knote@ht#1{\dimexpr\ht#1+\dp#1\relax}
\def\knote@setheight#1{\global\pdfpageheight=#1\relax
  \special{pdf:pagesize width \the\paperwidth\space height \the\pdfpageheight}}
\def\knote@flush{\ifvoid\knote@page\else
  \knote@setheight{\dimexpr\knote@ht\knote@page+40mm+2mm\relax}%
  \box\knote@page\par\clearpage\global\pdfpageheight=\paperheight\fi}
\newenvironment{knotechunk}
  {\par\setbox\knote@chunk=\vbox\bgroup\hsize=\textwidth\linewidth=\textwidth}
  {\par\egroup
   \ifdim\knote@ht\knote@chunk>\knote@limit
     \knote@flush\unvbox\knote@chunk\par\clearpage
   \else\ifvoid\knote@page
     \global\setbox\knote@page=\vbox{\unvbox\knote@chunk}%
   \else\ifdim\dimexpr\knote@ht\knote@page+\knote@ht\knote@chunk\relax>\knote@limit
     \knote@flush\global\setbox\knote@page=\vbox{\unvbox\knote@chunk}%
   \else
     \global\setbox\knote@page=\vbox{\unvbox\knote@page\unvbox\knote@chunk}%
   \fi\fi\fi}
\AtEndDocument{\knote@flush}
\makeatother
\pagestyle{empty}
\begin{document}
\begin{knotechunk}
{\color{knoteaccent}\rule{\linewidth}{1.5pt}}\par
{\LARGE\bfseries Mechanical properties — Dislocations and strengthening}\par
{\large Lecturer: Prof. J. Adeyemi}\hfill{\color{knotemuted}16 October 2026 \textperiodcentered{} Materials building, LT3}\par
{\color{knoteaccent}\rule{\linewidth}{0.6pt}}\par

\section{Stress and strain}

Engineering stress $\sigma = F/A_0$, strain $\varepsilon = \Delta L/L_0$; in the elastic range Hooke's law $\sigma = E\varepsilon$.

\begin{itemize}\setlength{\itemsep}{0pt}\setlength{\parskip}{0pt}
  \item yield strength: the 0.2 \% offset stress
  \item ultimate tensile strength (UTS): the maximum of the engineering curve
  \item ductility: elongation at fracture
\end{itemize}

\knotekey{Metals yield at stresses far below the theoretical shear strength (\ensuremath{\approx} G/10) — because of dislocations.}
\end{knotechunk}

\begin{knotechunk}
\section{Dislocations and slip}

\begin{itemize}\setlength{\itemsep}{0pt}\setlength{\parskip}{0pt}
  \item edge dislocation: Burgers vector $\mathbf{b} \perp$ line; screw: $\mathbf{b} \parallel$ line
  \item slip happens on close-packed planes along close-packed directions
  \item FCC: $\{111\}\langle 110\rangle$, 12 slip systems — hence ductile (Cu, Al, Au)
  \item Schmid's law: resolved shear stress $\tau_R = \sigma\cos\phi\cos\lambda$; slip starts when $\tau_R$ reaches the critical value
\end{itemize}
\end{knotechunk}

\begin{knotechunk}
\section{Strengthening mechanisms}

Every mechanism makes dislocations harder to move:

\begin{enumerate}\setlength{\itemsep}{0pt}\setlength{\parskip}{0pt}
  \item grain refinement — Hall–Petch: $\sigma_y = \sigma_0 + k_y d^{-1/2}$
  \item solid-solution strengthening: solute atoms strain the lattice
  \item precipitation hardening: particles pinned or bypassed by Orowan bowing (Al alloys, Ni superalloys)
  \item work hardening: dislocation density $\rho$ rises, $\tau \approx \alpha G b\sqrt{\rho}$ (Taylor)
\end{enumerate}

\knotekey{Most mechanisms trade strength for ductility; grain refinement improves both, down to the nanocrystalline range.}

\knotequestion{Why does Hall–Petch break down for grain sizes below about 10–20 nm?}
\end{knotechunk}
\end{document}
